

\chapter{保罗書信總結}
\section{保罗書信縱覽}


\subsection{寫作的時間順序}%教會書、罗马教會信}

%\begin{comment}
\iffalse
\def\rowi  {林前 & 所提尼 & 提摩太 &亞居拉百基拉&以弗所\\}
\def\rowii  {林後 &提摩太 & 提多及弟兄 &&馬其頓\\}
\def\rowviii {帖前 & 西拉提摩太 &&& \\}
\def\rowviiii{帖後 & 西拉提摩太 & && \\}
\def\rowiii{加拉太 & 保羅& &&\\}
\def\rowiv{以弗所 & 保羅& 推基古&&羅馬 \\}
\def\rowv{腓立比 & 提摩太& 提摩太以巴弗提&&羅馬\\}
\def\rowvi{羅馬書 & （德丟代） & 菲比&提摩太路求耶孫 &\\}
\def\rowvii{歌羅西 & 提摩太&推基古阿尼西母 &亞里達古馬可以巴弗猶士都&羅馬 \\}
\begin{table}[!hbp]
\begin{tabular}{lllll}
\toprule
書名 & 寫信人 & 送信人&同工&地點\\
\midrule
\sortfirstcol{\rowi\rowii\rowiii\rowiv\rowv\rowvi\rowvii\rowviii\rowviiii}
\bottomrule
\end{tabular}
\caption{教會書信}
\end{table}
%\end{comment}
\fi


\subsubsection{二次宣教及之前的書信(3)}
\vspace*{-\baselineskip}
\def\rowi  {帖前5 & 西拉，提摩太&二次宣教途中&哥林多&感恩，勸勉，復活及主再來  \\}
\def\rowii {帖後3 & 西拉，提摩太 &二次宣教途中&哥林多&離道反教，沉淪之子\\}
\def\rowiii{加拉太6 & 保羅 &一次宣教之後，耶路撒冷大會之前&敘利亞的安提阿&隨從別的福音\\}
\begin{table}[!hbp]
\begin{tabular}{lllll}
\toprule
書名 & 寫信人 & 寫作時間&寫作地點&主題\\
\midrule
\sortfirstcol{\rowi\rowii\rowiii}
\bottomrule
\end{tabular}
\caption{二次宣教及之前的書信}
\end{table}

\vspace*{-\baselineskip}
\subsubsection{三次宣教途中的書信(3)}
\vspace*{-\baselineskip}
\def\rowi  {林前16 & 所提尼 & 提摩太 &亞居拉，百基拉，司提反，福徒拿都，亞該古&以弗所&解決教會問題\\}
\def\rowii  {林後13 &提摩太 & 提多及弟兄 &&馬其頓&建立領袖同工\\}
\def\rowiii{羅馬書16 & （德丟代） & 菲比&提摩太，路求，耶孫，該猶，以拉都，括土， &哥林多&猶太人與外邦人\\}
\begin{table}[!hbp]
\begin{tabular}{llllll}
\toprule
書名 & 寫信人 & 送信人&同工&地點&主題\\
\midrule
\sortfirstcol{\rowi\rowii\rowiii}
\bottomrule
\end{tabular}
\caption{三次宣教途中的書信}
\end{table}

%\vspace*{-\baselineskip}
\subsubsection{監獄書——第一次在羅馬作監(4)}
\vspace*{-\baselineskip}
\def\rowi{腓立比4 & 提摩太& 提摩太，以巴弗提&&喜樂\\}
\def\rowii{以弗所6 & 保羅& 推基古&&榮耀，教會是基督的身體 \\}
\def\rowiii{歌羅西4 & 提摩太&推基古，阿尼西母 &亞里達古，馬可，以巴弗，猶士都，路加，底馬&豐盛，基督是教會的元首 \\}
\def\rowiv{腓利門書 & 提摩太 &推基古，阿尼西母&馬可,以巴弗,亞里達古,路加，底馬&\\}
\begin{table}[!hbp]
\begin{tabular}{lllll}
\toprule
書名 & 寫信人 & 送信人&同工&主題\\
\midrule
\sortfirstcol{\rowi\rowii\rowiii\rowiv}
\bottomrule
\end{tabular}
\caption{監獄書信}
\end{table}



\subsubsection{個人書信——第二次在羅馬作監(教牧書信)}
\begin{table}[H]
%\begin{tabularx}{12cm}{|l|l|l|l|l|X|X|}%{\textwidth}{|l|l|l|l|l|X|X|}%{12cm}{|X|X|X|X|X|X|X|}
\begin{tabularx}{17cm}{lllllll}%{\textwidth}{|l|l|l|l|l|X|X|}%{12cm}{|X|X|X|X|X|X|X|}
%\begin{tabularx}{12cm}{lllllXX}
\hline
書名 & 寫信人 &收信人&收信地點&寫作&送信人&同工\\ \hline
%腓 & 提摩太 &腓利門 &歌羅西& 羅馬&推基古，阿尼西母&馬可,以巴弗,亞里達古,路加，底馬\\%\hline
提前6 & 保羅 & 提摩太 &以弗所&馬其頓&&\\%\hline
提多3 & 保羅 &提多&克里特&馬其頓&推基古，或亞提馬& \\%\hline
提後4 & 保羅 & 提摩太 &以弗所&羅馬& 推基古& 路加,友布羅、布田、利奴、革老底亞\\\hline
\end{tabularx}
\caption{個人書信}
\end{table}

\subsection{解決的核心問題}




\subsubsection{羅馬書}
\begin{itemize}%[noitemsep]
\itemsep-.25em 
\item 猶太人的問題
\begin{itemize}%[noitemsep]
\itemsep-.25em 
\item 猶太人自認為有律法和割禮就論斷人 

\begin{itemize}%[noitemsep]
\itemsep-.25em 
%\item 自恃熟悉律法而论断人 2:1-11
\item 对有无律法的审判标准不清楚 2:12-16
\item 自恃熟悉律法却知法犯法 2:17-24
\item 未受割礼却全守律法的人要审判受过割礼却犯律法的人 2:25-27
\item 真割礼和真犹太人的定义 2:28-29
\end{itemize}

\item 因律法稱義還是靠信心 
\begin{itemize}%[noitemsep]
\itemsep-.25em
\item 犹太人的不信不能废掉神的信 3:1-4
\item 犹太人的不信反而显出神的荣耀，神发怒不合理（犹太人虽然不信，但他们守住了律法，所以神不应当因为他们不信而发怒） 3:5-8
\item 对犹太人有罪的概念不清楚 3:9-20 
\item 神在律法以外显明的义 3:9-20 
\item 因信耶稣而坚固律法 3:21-31
\end{itemize}
\end{itemize}

\item 外邦人的問題
\begin{itemize}%[noitemsep]
\itemsep-.25em 
\item 以色列人的不信說明神的話落空了嗎？
\item 神棄絕以色列人了嗎？ 
\end{itemize}
\end{itemize}

\subsubsection{以弗所書?}


\subsubsection{加拉太教會}
\begin{itemize}
\itemsep-.5em 
\item 這麽快離棄福音，隨從別的福音
\item 靠律法稱義而非信心
\item 守割禮和其他猶太律法的問題
\end{itemize}

\subsubsection{林前}%要解決的問題
\begin{itemize}%[noitemsep]
\itemsep-.25em 
\item 解決保羅從信徒口中聽到的哥林多的問題（1-6章）
\begin{itemize}
\itemsep-.25em 
\item 分黨-中間有紛爭。
\item 淫亂-有人收了他的繼母。
\item 告狀-彼此相爭，告在不信主的人面前
\end{itemize}
\item  解決哥林多教會信中的問題（7-16章）
\begin{itemize}
\itemsep-.25em 
\item 婚姻問題
\item 身份
\item 吃祭偶像之物
\item 女人蒙頭
\item 混亂聖餐
\item 屬靈恩賜
\item 死人復活
\item 給耶路撒冷教會捐錢的事
\end{itemize}
\end{itemize}

\subsubsection{林後}%要解決的問題
\begin{itemize}%[noitemsep]
\itemsep-.25em 
\item 保羅為自己的辯解（1-6章）
\begin{itemize}
\itemsep-.25em 
\item 兩次修改行程的原因
\item 自己成為新約的執事
\item 在各樣事上表明自己是神的用人
\end{itemize}
\item 勸哥林多樂意捐款給耶路撒冷教會（7-9章）
\item 警告哥林多自己要第三次再去哥林多（10-13章）
\begin{itemize}
\itemsep-.25em 
\item 保羅與假使徒的區別
\item 第三次要帶著使徒權柄再去哥林多
\end{itemize}
\end{itemize}
%\subsubsection{林后}




\subsubsection{帖前}%要解決的問題
\begin{itemize}%[noitemsep]
\itemsep-.25em 
\item 可能對保羅快速的離開他們有誤解“以爲保羅怕死”？
\item 成為聖潔，遠避淫行，不要沾染汙穢
\item 因不明白復活而為睡了的人而憂傷
\item 勸他們驚醒，等候主的再來
\end{itemize}


\subsubsection{帖后}%要解決的問題
\begin{itemize}%[noitemsep]
\itemsep-.25em 
\item 在患難和逼迫中要站立得穩
\item 冒充保羅的名口傳或者寫信說耶穌的日子已經到了
\item 有人不做工，反倒專管閑事
%\item  4 勸他們驚醒，等候主的再來
\end{itemize}


\subsubsection{腓立比教會}%的問題
\begin{itemize}%[noitemsep]
\itemsep-.25em 
\item 保羅坐監是為使福音興旺，鼓勵腓立比人不要懼怕
\item 同心合意不發怨言
\item 不要靠肉體和律法稱義，而靠信心
\item 勸兩個女人尤阿迪婭和循都基要在主裡同心
\end{itemize}

\subsubsection{歌羅西書}%以弗所書要解決的問題}
\begin{itemize}%[noitemsep]
\itemsep-.25em 
\item 有人傳理學和虛空的妄言，不照著基督，乃照人間的遺傳(西2:8）
\begin{itemize}%[noitemsep]
\itemsep-.25em 
\item 注重割禮(西2:11）
\item 注重飲食，或節期、月朔、安息日(西2:16）
\item 故意謙虛和敬拜天使，。這等人拘泥在所見過的，隨著自己的慾心，無故的自高自大，(西2:18） 
\item 服從那「不可拿、不可嘗、不可摸」等類的從人而來的規條(西2:20）
\end{itemize}
\item 信徒生活方面
\begin{itemize}%[noitemsep]
\itemsep-.25em 
\item 淫亂、污穢、邪情、惡慾，和貪婪（貪婪就與拜偶像一樣）(西3:5）
\item 心中存惱恨、忿怒、惡毒、毀謗，並口中污穢的言語，彼此說謊(西3:8,9）
\end{itemize}
\end{itemize}

\subsubsection{腓利門書}
保羅盼望腓利門接納逃奴阿尼西姆。

\subsubsection{提前——解决教会领袖的问题}
\begin{itemize}%[noitemsep]
\itemsep-.25em 
\item 幾個人傳異教，講荒渺無憑的話語和無窮的家譜,虛浮的話 
\item 有人講虛浮的話,想要作教法師，卻不明白自己所講說的所論定的
\item 許米乃和亞力山大等人丟棄良心，就在真道上如同船破壞了一般 
\item 挑選監督和執事的資格 
\item 如何教導勸解教會裡不同類型的人
\item 勸提摩太如何自己做信徒的榜樣
\end{itemize}

\subsubsection{提多——解决教会会友的问题}
\begin{itemize}%[noitemsep]
\itemsep-.25em 
\item 在各城設立長老 
\item 監督和長老的資格
\item 如何在教會勸解各種類型的信徒
\end{itemize}

\subsubsection{提後——教会领袖的榜样}
\begin{itemize}%[noitemsep]
\itemsep-.25em 
\item 勸提摩太剛強壯膽，不要膽怯,不以福音為恥
\item 守住善道，做基督的精兵 
\item 末世教會光景
\end{itemize}


